\section{
    Almost simply typed almost reflective language:
    \texorpdfstring{$\SSTLC$}{[[S]]TLC}
}\label{SSTLC}

We start with a simpler language to iron out the details
and actually figure out how to properly reinterpret
terms internally before moving on to a dependently typed
system. The simplest language for which we know a lot of
interesting models is the simply typed lambda calculus
with products, which we refer to as ``STLC'' throughout
this paper; in this section, we'll present a language
capable of interpreting STLC terms as constructions in
STLC models that can be defined in our language.

So, what are the models of STLC? In addition to the
usual set-theoretic semantics where $S\to T$ is
interpreted as a function from $S$ to $T$ and
$S\times T$ is a cartesian product, there are
proof-theoretic semantics that relate terms of type $T$
to intuitionistic proofs of proposition
$\llbracket T\rrbracket$, where $(\to)$ is now
understood to mean implication, and $(\times)$ is
conjunction. This, of course, is a classic Curry-Howard
correspondence \cite{ho80}; both set-theoretic and
proof-theoretic semantics are generalized by a complete
description of STLC models as given by Joachim Lambek
\cite{Lambek1980}. According to this
\textit{Curry-Howard-Lambek} correspondence, models
of STLC are precisely the cartesian closed categories
(or CCCs for short). Definitions vary; we're going to
use the definition \ref{ccc}.

\begin{definition}[Cartesian closed category]\label{ccc}
A cartesian closed category is the following datum:
\begin{enumerate}
\item A class\footnote{Objects do not necessarily form a
    set: for example, one might want to have a category
    of sets where morphisms are functions between them.
    To deal with sizing issues, pick appropriate
    set-theoretical foundations you are most comfortable
    with (NBG, MK, ZFC/S or something else); this detail
    is not important in our development} of
    \textit{objects} $\mathrm{Ob}$ and a family of
    \textit{morphisms} $\mathrm{Hom}(S, T)$ indexed by
    $S,T\in\mathrm{Ob}$, where $f\in\mathrm{Hom}(S,T)$
    is denoted by $f:S\to T$. In addition, $S$ is called
    the \textit{source} of $f$, and $T$ is the
    \textit{target} of $f$.
\item For each $A\in\mathrm{Ob}$ there's an
    \textit{identity morphism} $\mathrm{id}_A: A\to A$;
    for each $f:A\to B$ and $g:B\to C$, there's their
    \textit{composition} $g\circ f:A\to C$. If a
    source-target of an identity morphism is clear from
    the context, the subscript is elided. Also, the
    following properties hold:
    \begin{align*}
        \mathrm{id}\circ f = f = f\circ\mathrm{id} & &
        (f \circ g) \circ h = f \circ (g \circ h)
    \end{align*}
    This applies whenever all compositions are well-defined.
\item There's a \textit{terminal object}
    $\top\in\mathrm{Ob}$ such that for each
    $A\in\mathrm{Ob}$ there's a \textbf{unique} morphism
    $!_A:A\to\top$.
\item For each $A,B\in\mathrm{Ob}$ there's their
    \textit{product} $A\times B$ together with
    \textit{projections} $\pi_1:A\times B\to A$ and
    $\pi_2:A\times B\to B$; for each object
    $C\in\mathrm{Ob}$ and a pair of morphisms $f:C\to A$
    and $g:C\to B$, there's a \textbf{unique}
    \textit{pairing morphism} $[f;g]:C\to A\times B$
    such that the following properties hold:
    \begin{align*}
        \pi_1\circ[f;g]=f & & \pi_2\circ[f;g]=g
    \end{align*}
\item For each $A,B\in\mathrm{Ob}$ there's their
    \textit{exponential object} $B^A$ together with
    \textit{evaluation morphism}
    $\mathrm{ev}:B^A\times A\to B$; for each object
    $C\in\mathrm{Ob}$ and morphism $f:C\times A\to B$,
    there's a \textbf{unique} \textit{abstraction}
    $\lambda\,f:C\to B^A$ such that the following
    property holds for arbitrary $g:C\to A$:
    \[
        \mathrm{ev}\circ[\lambda\,f; g]
        = f\circ[\mathrm{id}_C; g]
    \]
\end{enumerate}
\end{definition}

Then, a well-typed term $\Gamma\vdash t:T$ is
interpreted as a morphism $\llbracket t\rrbracket:
\llbracket\Gamma\rrbracket\to\llbracket T\rrbracket$,
where $\llbracket\Gamma\rrbracket$ is a left-nested
product of interpretations of types from $\Gamma$:
\[
    \llbracket x:A,y:B,z:C\rrbracket=
    ((\top\times\llbracket A\rrbracket)
    \times\llbracket B\rrbracket)
    \times\llbracket C\rrbracket.
\]
Unit type is now a terminal object; a product of types is
now a product of objects; an arrow type is an exponential
object. Equations in (2) are needed to validate the
structural properties of substitution; equations in (4)
correspond to reductions of pairs; equation in (5)
validates $\beta$-reduction. \textbf{Unique}ness
properties correspond to $\eta$-equivalence.

For example, in our language, $m\,\llbracket t
\rrbracket$ is the operator interpreting the term $t$ in
an internal model $m$. On a sample term
\[
    \cdot\vdash\lambda f x y. (f x, y):
        (A\to B)\to A\to C\to B\times C
\]
it computes like this:
\begin{multline*}
    m\,\llbracket\lambda f x y. (f x, y))\rrbracket=\\=
    \mathtt{lam}\,(\mathtt{lam}\,(\mathtt{lam}\,(
        \mathtt{pair}\,
            (\mathtt{eval}\,\circ\,\mathtt{pair}\,
                (\mathtt{proj}_2\circ\mathtt{proj}_1
                    \circ\mathtt{proj}_1
                )\,
                (\mathtt{proj}_2\circ\mathtt{proj}_1)
            )\,
            \mathtt{proj}_2
    )))
\end{multline*}

Where $\mathtt{lam}$,\,$\mathtt{pair}$ and
$\mathtt{eval}$ are the abstraction morphism, the
pairing morphism, and the evaluation morphism taken from
$m$, defined in the same way as in the definition
\ref{ccc}, translated into types of our system.

Now, moving on to the actual design of our language. As
discussed in Section~\ref{no-go-section}, for the
non-dependent case we are now dealing with, we will
interpret only normal forms of closed terms. Among other
things, this means that we actually don't need equations
and \textbf{unique}nesses from Definition~\ref{ccc} to
hold in our internal models! This approach doesn't break
any of the metatheoretical properties considered in
Section~\ref{SSTLC-properties}. Additionally, for the
sake of example, we're only going to consider models
where objects are types. Then, the internalization of
definition \ref{ccc} would require 1) type operators to
define types of morphisms, products, and exponential
objects; 2) polymorphic functions to operationalize the
$\forall\exists$ introductions of constructions in a
model. Actually, we're going to get rid of type
operators as well, since we can do with just
types-in-a-context which are then substituted into
the definition of a model for checking term-level
operations. This way, our language is a small extension
of the classic Hindley-Milner type system; grammar rules
are collected in a Figure~\ref{sstlc-grammar}. You'll
notice that, despite our initial description, grammar
doesn't define a reinterpretation operator
$m\llbracket t\rrbracket$ for arbitrary term $t$;
instead, only variables (denoted by a metavariable $x$)
get to be reinterpreted. In this system,
reinterpretation is nothing but a total function on
closed values which is called whenever a value is
substituted into such an $x$ (so, in a way, this is
similar to hereditary substitution
\cite{watkins2003concurrent,keller2010hereditary}).

\begin{definition}[Values of $\SSTLC$]
Values are denoted by metavariables $u,v$ in the BNF
grammar below. Metavariable $n$ denotes neutral terms:
\begin{align*}
    u,v&\Coloneqq n\;|\;()\;|\;(u,v)\;|\;\lambda x.v\\
    n&\Coloneqq x\;|\;\pi_i\,n\;|\;n\,v
\end{align*}
\end{definition}

\begin{figure}
\centering
\begin{bnf}
    $S,T,C,P,E,TT$ : \textsf{MonoType} ::=
    | $\alpha$ : type variable
    | $\top$ : unit type
    | $S\times T$ : product type
    | $S\to T$ : arrow type
    ;;
    $\Delta$ : \textsf{MonoCtx} ::=
    | $\cdot$ : empty context
    | $\Delta,x\colon T$ : context extension
    ;;
    $\tau$ : \textsf{PolyType} ::=
    | $T$ : monomorphic type
    | $\forall\alpha.\tau$ : forall type
    ;;
    $\Gamma$ : \textsf{PolyCtx} ::=
    | $\cdot$ : empty context
    | $\Gamma,x\colon\tau$ : context extension
    ;;
    $\Xi$ : \textsf{Ctx} ::=
    | $\Gamma\,\|\,\Delta$ : joined context
    ;;
    $s,t$ : \textsf{Term} ::=
    | $x$ : variable
    | $()$ : singleton
    | $(s, t)$ : pair
    | $\pi_1\,t$ : first projection
    | $\pi_2\,t$ : second projection
    | $\lambda x. t$ : abstraction
    | $t\,s$ : application
    | \texttt{let}\,$x=s$\,\texttt{in}\,$t$
        : polymorphic binding
    | $m\llbracket x\rrbracket$ : reinterpretation
    ;;
    $m$ : \textsf{Model} ::=
    | $\langle T, t, \ldots \rangle$
        : model def. (see typing rules)
\end{bnf}
\caption{$\SSTLC$ terms and values}
\label{sstlc-grammar}
\end{figure}

Thus, to define interpretation, it is enough to consider
STLC terms:

\begin{definition}[Interpretation of terms in $\SSTLC$]
    Given a model $m$, a list of variables $\gamma$ and
    a term $t$, interpretation of a term $m\llbracket
    \gamma\vdash t\rrbracket$ is a recursive function
    defined by the following equations:
    \begin{align}{}
        m\llbracket\gamma\vdash()\rrbracket&=
        \mathtt{cnst}
        \\
        m\llbracket\gamma\vdash(s,t)\rrbracket&=
        \mathtt{pair}
        \,m\llbracket\gamma\vdash s\rrbracket
        \,m\llbracket\gamma\vdash t\rrbracket
        \\
        m\llbracket\gamma\vdash\pi_i\,t\rrbracket&=
        \mathtt{comp}\,\mathtt{proj}_i\,
        m\llbracket\gamma\vdash t\rrbracket
        \\
        m\llbracket\gamma\vdash x\rrbracket&=\gamma[x]
        \\
        m\llbracket\gamma\vdash\lambda x.t\rrbracket&=
        \mathtt{lam}\,
        m\llbracket(\gamma,x)\vdash t\rrbracket
        \\
        m\llbracket\gamma\vdash t\,s\rrbracket&=
        \mathtt{comp}\,\mathtt{eval}\,\left(
        \mathtt{pair}
        \,m\llbracket\gamma\vdash t\rrbracket
        \,m\llbracket\gamma\vdash s\rrbracket
        \right)
    \end{align}
    where model operations \texttt{cnst},\,\texttt{pair}
    and others are taken from $m$ and $\gamma[x]$ is
    also defined by equations:
    \begin{align}
        (\gamma,x)[x] &= \mathtt{proj}_2 \\
        (\gamma,y)[x] &= \mathtt{comp}\,\gamma[x]\,
        \mathtt{proj}_1
    \end{align}
\end{definition}

Here is our way to enforce a principle that \textit{only
closed terms get interpreted}: once a term $s$ bound by
a polymorphic let-binding
$\mathtt{let}\,x=s\,\mathtt{in}\,t$ becomes closed and
is brought into normal form, only then is it
substituted into $t$, and whenever it is substituted
into $m\llbracket x\rrbracket$, it immediately gets
reinterpreted. This operation is, as noted before,
similar to hereditary substitution; it is denoted by
$t[s\sslash x]$ in our computation rules given in
Figure~\ref{sstlc-comp}. So it acts like a regular
substitution in all cases except for variables under
reinterpretation:
\[
    m\llbracket x\rrbracket[s\sslash x] =
    m\llbracket()\vdash s\rrbracket
\]
One can easily notice that computation would get stuck
if let-bound terms were to contain variables coming
from unapplied $\lambda$-abstractions. Our typing rules
given in Figure~\ref{sstlc-typing} protect us from that,
allowing to refer only to let-bound terms inside other
let-bound terms. Again, aside from that and the typing
rule for universes, there's nothing different from the
usual declarative typing rules for a Hindley-Milner type
system \cite{Clement:ML1986,damas1982principal,
milner1978theory}.

\begin{figure}
\centering
\begin{mathpar}
    \inferrule{s\leadsto s'}{(s,t)\leadsto(s',t)}
    \and
    \inferrule{t\leadsto t'}{(s,t)\leadsto(s,t')}
    \and
    \inferrule{t\leadsto t'}{\pi_i t\leadsto\pi_i t'}
    \and
    \inferrule{}{\pi_i(t_1, t_2)\leadsto t_i}
    \\
    \inferrule{t\leadsto t'}
        {\lambda x.t\leadsto\lambda x.t'}
    \and
    \inferrule{s\leadsto s'}{s\,t\leadsto s'\,t}
    \and
    \inferrule{t\leadsto t'}{s\,t\leadsto s\,t'}
    \and
    \inferrule{}{(\lambda x.t)s\leadsto t[s/x]}
    \\
    \inferrule{s\leadsto s'}{
        \mathtt{let}\,x=s\,\mathtt{in}\,t\leadsto
        \mathtt{let}\,x=s'\,\mathtt{in}\,t
    }
    \and
    \inferrule{t\leadsto t'}{
        \mathtt{let}\,x=s\,\mathtt{in}\,t\leadsto
        \mathtt{let}\,x=s\,\mathtt{in}\,t'
    }
    \and
    \inferrule{\mathrm{fv}(v)=\varnothing}{
        \mathtt{let}\,x=v\,\mathtt{in}\,t
        \leadsto t[v\sslash x]
    }
\end{mathpar}
\caption{$\SSTLC$ computation rules}
\label{sstlc-comp}
\end{figure}

To check models, we abuse our type system a bit 
to avoid adding type operators and complicating it:
types of morphisms $C$, product objects $P$ and
exponentials $E$ are specified as open monomorphic types
which can use two type variables, $0$ and $1$. This fact
is reflected in the syntax of a model using notation
$\Lambda 01.C$, $\Lambda 01.P$ and $\Lambda 01.E$
respectively. $0$ and $1$ are then filled in by
substitution when we need to type model operations
(notation $T[S_1, S_2]$ is a shorthand for
$T[S_1/0, S_2/1]$). The terminal object is
just a type named $TT$.

In order to check reinterpretation, we first need to
define how the model acts on monomorphic types. As there
is no computation on them, reinterpretation of types is
defined once again as a recursive function on syntax.

\begin{definition}[Interpretation of types in $\SSTLC$]
    Let $m$ be a correct model and $\sigma$ be a
    substitution of monomorphic types. Then, the
    interpretation of a type in a model is the following
    function:
    \begin{align}{}
        m\llbracket\sigma\vdash\alpha\rrbracket&=
            \sigma\alpha \\
        m\llbracket\sigma\vdash\top\rrbracket&=TT \\
        m\llbracket\sigma\vdash(S\times T)\rrbracket&=
            P[m\llbracket\sigma\vdash S\rrbracket,
            m\llbracket\sigma\vdash T\rrbracket] \\
        m\llbracket\sigma\vdash S\to T\rrbracket&=
            E[m\llbracket\sigma\vdash S\rrbracket,
            m\llbracket\sigma\vdash T\rrbracket].
    \end{align}
\end{definition}

\begin{figure}
\centering
\begin{mathpar}
    \inferrule{}{\Xi\vdash ():\top}
    \and
    \inferrule
    {\Xi\vdash s:S\\\Xi\vdash t:T}
    {\Xi\vdash (s,t):S\times T}
    \and
    \inferrule
    {\Xi\vdash t:S\times T}
    {\Xi\vdash \pi_1\,t:S}
    \and
    \inferrule
    {\Xi\vdash t:S\times T}
    {\Xi\vdash \pi_2\,t:T}\\
    \and
    \inferrule{x:T\in\Delta}
    {\Gamma\,\|\,\Delta\vdash x:T}
    \and
    \inferrule{\Gamma\,\|\,\Delta,x:S\vdash t:T}
        {\Gamma\,\|\,\Delta\vdash\lambda x.t:S\to T}
    \and
    \inferrule
    {\Xi\vdash t:S\to T\\\Xi\vdash s:S}
    {\Xi\vdash t\,s:T}
    \and
    \inferrule{
        \Gamma\vdash s:S\\
        \mathrm{fv}(S)=\vec{\alpha}\\
        \Gamma,x:\forall\vec{\alpha}.S\,\|\,\Delta
        \vdash t:T
    }{\Gamma\,\|\,\Delta\vdash
        \mathtt{let}\,x=s\,\mathtt{in}\,t:T}
    \and
    \inferrule{
        \Gamma\,\|\,\Delta\vdash m\,\mathrm{model}\\
        x:\forall\vec{\alpha}.T\in\Gamma\\
        \sigma=\left[\vec{S}/\vec{\alpha}\right]
    }{\Gamma\,\|\,\Delta\vdash m\llbracket x\rrbracket
        :C_m[TT_m, m\llbracket\sigma\vdash T\rrbracket]}
    \and
    \inferrule{
        \mathrm{fv}(C)\cup\mathrm{fv}(P)\cup
            \mathrm{fv}(E)\subseteq\{0,1\}\\
        \Xi\vdash\mathtt{id}:C[\alpha,\alpha]\\
        \Xi\vdash\mathtt{comp}:
            C[\beta,\gamma]\to C[\alpha,\beta]\to
            C[\alpha,\gamma]\\
        \mathrm{fv}(TT)=\varnothing\\
        \Xi\vdash\mathtt{cnst}:C[\alpha,TT]\\
        \Xi\vdash\mathtt{pair}:
            C[\gamma,\alpha]\to C[\gamma,\beta]\to
            C[\gamma,P[\alpha,\beta]]\\
        \Xi\vdash\mathtt{proj}_1:
            C[P[\alpha,\beta],\alpha]\\
        \Xi\vdash\mathtt{proj}_2:
            C[P[\alpha,\beta],\beta]\\
        \Xi\vdash\mathtt{lam}:
            C[P[\gamma,\beta],\alpha]\to
            C[\gamma,E[\alpha,\beta]]\\
        \Xi\vdash\mathtt{eval}:
            C[P[E[\alpha,\beta],\beta],\alpha]
    }{\Xi\vdash\langle
        \Lambda 01.C,\mathtt{id},\mathtt{comp}, TT,
        \mathtt{cnst}, \Lambda 01.P,\mathtt{pair},
        \mathtt{proj}_1,\mathtt{proj}_2,
        \Lambda 01.E,\mathtt{lam},\mathtt{eval}
    \rangle\,\mathrm{model}}
\end{mathpar}
\caption{$\SSTLC$ declarative typing rules}
\label{sstlc-typing}
\end{figure}

\subsection{Metatheoretic properties of
\texorpdfstring{$\SSTLC$}{[[S]]TLC}}
\label{SSTLC-properties}

Here we only study some basic properties of a system to
make sure that everything works as expected:
computations are confluent; types are preserved by all
computations; all the types can be inferred if you
wish to do so; everything terminates, and, modulo
reduction of \texttt{let}-bindings, this is the same old
STLC we started with. Proofs mostly follow
standard techniques since the rules for reinterpretation
are structured to directly plug into proofs of
properties of the STLC and the Hindley-Milner type
system.

\begin{theorem}[Church-Rosser theorem for $\SSTLC$]
    Reflexive transitive closure of $(\leadsto)$ has
    the Church-Rosser property.
\end{theorem}
\begin{proof}
    Classic proof via parallel reduction relation
    \cite{takahashi1995parallel,barendregt1984lambda}.
    Parallel reduction is defined as follows:
    \begin{mathpar}
        \inferrule{}{t\rightrightarrows t}
        \and
        \inferrule{
            s\rightrightarrows s'\\
            t\rightrightarrows t'
        }{(s,t)\rightrightarrows(s',t')}
        \and
        \inferrule{t\rightrightarrows t'}
            {\pi_i\,t\rightrightarrows\pi_i\,t'}
        \and
        \inferrule{t_1\rightrightarrows t_1'}
            {\pi_1(t_1,t_2)\rightrightarrows t_1'}
        \and
        \inferrule{t_2\rightrightarrows t_2'}
            {\pi_2(t_1,t_2)\rightrightarrows t_2'}
        \\
        \inferrule{t\rightrightarrows t'}
            {\lambda x.t\rightrightarrows\lambda x.t'}
        \and
        \inferrule{
            s\rightrightarrows s'\\
            t\rightrightarrows t'
        }
        {s\,t\rightrightarrows s'\,t'}
        \and
        \inferrule{
            s\rightrightarrows s'\\
            t\rightrightarrows t'
        }
        {(\lambda x.s)\,t\rightrightarrows
            s'[t'/x]}
        \\
        \inferrule{
            s\rightrightarrows s'\\
            t\rightrightarrows t'
        }{
            \mathrm{let}\,x=s\,\mathrm{in}\,t
                \rightrightarrows
            \mathrm{let}\,x=s'\,\mathrm{in}\,t'
        }
        \and
        \inferrule{t\rightrightarrows t'}{
            \mathrm{let}\,x=v\,\mathrm{in}\,t
                \rightrightarrows
            t'[v\sslash x]
        }
    \end{mathpar}
\end{proof}

\begin{theorem}[Type preservation for $\SSTLC$]
    The following is true:
    \begin{enumerate}
        \item If $\Xi\vdash m\,\mathrm{model}$,
            $\Delta\vdash t:T$ and $\delta$ is a list of
            variables from $\Delta$, then \\ $\Xi\vdash
            m\llbracket\delta\vdash t\rrbracket:
            C_m[m\llbracket\Delta\rrbracket,
            m\llbracket T\rrbracket]$.
        \item If $\Gamma\,\|\,\Delta\vdash s:S$ and
            $\Gamma\,\|\,\Delta,x:S,\Delta'\vdash t:T$,
            then $\Gamma\,\|\,\Delta,\Delta'\vdash
            t[s/x]:T$.
        \item If $\Gamma\vdash v:\tau$ and
            $\Gamma,x:\tau,\Gamma'\,\|\,\Delta\vdash
            t:T$, then
            $\Gamma,\Gamma'\,\|\,\Delta\vdash
            t[v\sslash x]:T$.
        \item If $\Gamma\vdash t:T$ and $t\leadsto t'$,
            then $\Gamma\vdash t':T$.
    \end{enumerate}
\end{theorem}
\begin{proof}
    By structural induction on terms and $(\leadsto)$.
\end{proof}

\begin{theorem}[Strong normalization for $\SSTLC$]
    Closed terms typable in $\SSTLC$ are strongly
    normalizing with respect to $(\leadsto)$ and reduce
    to a value.
\end{theorem}
\begin{proof}
    Classic proof via saturated sets
    \cite{tait1967intensional,girard1989proofs,
    luo1994computation}. Let $\mathsf{SN}$ be a set of
    terms which are strongly normalizing with respect to
    $(\leadsto)$ and normal forms of which are values.
    Let $\mathsf{Atomic}$ be the least set of terms
    satisfying the following conditions: variables $x$
    are $\mathsf{Atomic}$; if $t$ is $\mathsf{Atomic}$,
    then so are $\pi_i\,t$ and $t\,s$. Then, let
    $t\to_\mathrm{WHR}t'$ be the following relation on
    terms:
    \begin{mathpar}
        \inferrule{t\in\mathsf{SN}}
        {(\lambda x.s)t\to_\mathrm{WHR}s[t/x]}
        \and\inferrule{t_2\in\mathsf{SN}}
        {\pi_1(t_1,t_2)\to_\mathrm{WHR}t_1}
        \and\inferrule{t_1\in\mathsf{SN}}
        {\pi_2(t_1,t_2)\to_\mathrm{WHR}t_2}
        \and\inferrule{t\to_\mathrm{WHR}t'}
        {t\,s\to_\mathrm{WHR}t'\,s}
        \and\inferrule{t\to_\mathrm{WHR}t'}
        {\pi_i\,t\to_\mathrm{WHR}\pi_i\,t'}
        \and\inferrule{ }
        {\mathtt{let}\,x=v\,\mathtt{in}\,t
        \to_\mathrm{WHR}t[v\sslash x]}
        (\mathrm{fv}(v)=\varnothing)
    \end{mathpar}

    Finally, let $\mathsf{SAT}$ be the following family
    of sets of terms:
    \begin{itemize}
        \item if $X\in\mathsf{SAT}$, then
            $X\subseteq\mathsf{SN}$;
        \item if $X\in\mathsf{SAT}$, then
            $\mathsf{Atomic}\cap\mathsf{SN}\subseteq
            \mathsf{SAT}$;
        \item if $X\in\mathsf{SAT}$, $t'\in X$ and
            $t\to_\mathrm{WHR}t'$, then $t\in X$.
    \end{itemize}

    Let there be the following relation
    $\mathcal{R}_\tau^\zeta\in\mathsf{SAT}$, where
    $\zeta$ is a mapping from type variables into
    $\mathsf{SAT}$:
    \begin{itemize}
        \item $t\in\mathcal{R}_\alpha^\zeta$ iff
            $t\in\zeta(\alpha)$;
        \item $t\in\mathcal{R}_\top^\zeta$ iff
            $t\in\mathsf{SN}$;
        \item $t\in\mathcal{R}_{T_1\times T_2}^\zeta$
            iff $t\in\mathsf{SN}$ and
            $\pi_i\,t\in\mathcal{R}_{T_i}^\zeta$
            for $i=1,2$;
        \item $t\in\mathcal{R}_{S\to T}^\zeta$ iff
            $t\in\mathsf{SN}$ and for all
            $s\in\mathcal{R}_S^\zeta$,
            $t\,s\in\mathcal{R}_T^\zeta$;
        \item $t\in\mathcal{R}_{\forall\alpha.T
            }^\zeta$ iff $t\in\mathsf{SN}$ and
            $t\in\mathcal{R}_T^{\zeta,\alpha\mapsto X}$
            for all $X\in\mathsf{SAT}$.
    \end{itemize}

    Indeed, $\mathcal{R}_\tau^\zeta\in\mathsf{SAT}$.
    Furthermore, the following property holds: if
    $t\in\mathcal{R}_{\forall\vec{\alpha}.T}^\zeta$,
    then $m\llbracket()\vdash\mathrm{nf}(t)\rrbracket
    \in\mathcal{R}^\zeta_{m\llbracket[\vec{S}/
    \vec{\alpha}]\vdash T\rrbracket}$ for all
    $m$ s.t. all terms inside it are in corresponding
    $\mathcal{R}_\forall^\zeta$.

    Then we can prove our main lemma: if
    $\Gamma\vdash t:T$ and $\gamma$ is a substitution
    s.t. $\gamma\,x\in\mathcal{R}_{\Gamma(x)}^\zeta$ for
    all $x\in\Gamma$, then $\gamma\,t\in
    \mathcal{R}_T^\zeta$.
\end{proof}

\begin{theorem}[Conservativity over STLC]
    Normal forms of closed terms typable in $\SSTLC$ are
    typable in STLC.
\end{theorem}
\begin{proof}
    Thanks to type preservation, normal form
    $\mathrm{nf}(t)$ is typable in $\SSTLC$. By
    definition, it is a value; therefore, it has no
    reinterpretation subterms and let-bindings, hence is
    typable in STLC.
\end{proof}

\begin{theorem}
    Global type inference for $\SSTLC$ is decidable.
\end{theorem}
\begin{proof}
    It is enough to extend any unification-based type
    inference algorithm for the Hindley-Milner type
    system to support models and reinterpretation.
    \begin{itemize}
        \item Since type operators in models are
            provided explicitly, it is enough to
            \textit{check} that their fields are typable
            with expected types.
        \item Since polymorphic functions do not depend
            on the local context, their type is fully
            resolved when they are added to the context.
        \item For case $m\llbracket x\rrbracket$, and
            $(x:\forall\vec{\alpha}.T)\in\Gamma$, we can
            directly compute $m\llbracket[\vec{\eta}/
            \vec{\alpha}]\vdash T\rrbracket$ (since
            $m$ is known and $T$ is fully resolved)
            where $\vec{\eta}$ are fresh metavariables.
    \end{itemize}
\end{proof}

Thus, in this section, we have introduced a sound type
system based on the Hindley-Milner type system which is
capable of interpreting its own closed terms in an
arbitrary ``cartesian closed category'' definable inside
it. This system is not truly \textit{reflective} since
models are not fully internal; they require a special
form of type operators to work, which is not accessible
to the rest of the type system. To internalize model
definition, we move on to a dependently typed system in
the next section.
